% !TEX TS-program = pdflatex
% !TeX program = pdflatex
% !TEX encoding = UTF-8
% !TeX spellcheck = en_US

\documentclass[xcolor=table]{beamer}

\usepackage{../extra/beamer/karimml}

\input{options2}

\subtitle[RNNs]{Recurrent Neural Networks}  

\changegraphpath{../img/RNN/}

\begin{document}

\begin{frame}
	\frametitle{\inserttitle}
	\framesubtitle{\insertshortsubtitle: Motivation}
	
	\begin{exampleblock}{Many AI tasks involve sequences}
		\small\vskip-8pt
		\begin{columns}
			\column{0.4\textwidth}
			\begin{itemize}
				\item Natural language:
				\begin{itemize}\footnotesize
					\item Machine translation
					\item Speech recognition
					\item Sentiment analysis
				\end{itemize}
			\end{itemize}
			
			\column{0.3\textwidth}
			\begin{itemize}
				\item Time series:
				\begin{itemize}\footnotesize
					\item Stock prices
					\item Weather forecasting
					\item Sensor signals
				\end{itemize}
			\end{itemize}
			
			\column{0.28\textwidth}
			\begin{itemize}
				\item Biological sequences:
				\begin{itemize}\footnotesize
					\item DNA
					\item Protein sequences
				\end{itemize}
			\end{itemize}
		\end{columns}
	\end{exampleblock}
	
	\vskip-6pt
	
	\begin{alertblock}{Problem}
		Traditional feed-forward neural networks process each input independently and cannot naturally model temporal dependencies or context.
	\end{alertblock}
	
	\vskip-6pt
	
	\begin{block}{Idea behind RNNs}
		Recurrent Neural Networks introduce a \textbf{memory mechanism} allowing information from previous time steps to influence current predictions.
	\end{block}
	
\end{frame}

\begin{frame}
	\frametitle{\inserttitle}
	\framesubtitle{\insertshortsubtitle: Definitions}
	
	\begin{block}{Definition}
		\vspace{6pt}
		A ``\textit{Recurrent Neural Network}'' is an architecture of Neural Networks. 
		In this one, there is a ``\textit{feed}'' back from the outputs of the neurons to the inputs throughout the network. \cite{2006-sazli}
		
		\vspace{6pt}
		In other word, there is a temporal phenomena treated in the problem in question; current output depends on the past one.
		\vspace{6pt}
	\end{block}
	
	\begin{exampleblock}{Example}
		\begin{itemize}
			\item Predicting future weather based on some features, plus the current one
		\end{itemize}
	\end{exampleblock}
	
\end{frame}

\begin{frame}
	\frametitle{\inserttitle}
	\framesubtitle{\insertshortsubtitle: Plan}
	
	\begin{multicols}{2}
		%	\small
		\tableofcontents
	\end{multicols}
\end{frame}

\section{Traditional sequence processing}

\begin{frame}
	\frametitle{\insertshortsubtitle}
	\framesubtitle{\insertsection}
	
	\begin{center}
		\hgraphpage[0.7\textwidth]{temporal_dependencies.pdf}
	\end{center}
	
%	\vspace{0.2cm}
	
	\begin{alertblock}{Key idea}
		The interpretation of the current element depends on previous observations.
	\end{alertblock}
	
\end{frame}

\subsection{Sequential and temporal data}

\begin{frame}
	\frametitle{\insertshortsubtitle: \insertsection}
	\framesubtitle{\insertsubsection}
	
	\begin{block}{Sequential data}
		Data where order matters. Observations are not independent.
	\end{block}
	
	\begin{itemize}
		
		\item \textbf{Sequence representation}
		\begin{itemize}
			\item Input shape: $(m, t, d)$ (batch × time steps × feature dimension)
			
			\item Output shape depends on the task:
			\begin{itemize}
				
				\item Sequence labeling: $(m, t, l)$ (batch × time steps × number of classes)
				
				\item Sequence classification: $(m, l)$ (batch × number of classes)
				
				\item Sequence generation: $(m, t', l)$ (batch × generated time steps × vocabulary/classes)
				
			\end{itemize}
		\end{itemize}
		
		\item \textbf{Variable-length sequences}
		\begin{itemize}
			\item Different samples may contain different numbers of time steps.
			\item Padding or masking is often used in mini-batch training.
		\end{itemize}
		
	\end{itemize}
	
\end{frame}

\begin{frame}
	\frametitle{\insertshortsubtitle: \insertsection}
	\framesubtitle{\insertsubsection: Contextual dependencies}
	
	\begin{exampleblock}{Example}
		``The animal did not cross the street because \textbf{it} was too tired.''
	\end{exampleblock}
	
	\vspace{0.3cm}
	
	The interpretation of ``it'' depends on previous words.
	
	\vspace{0.3cm}
	
	\begin{center}
		Current prediction $\leftarrow$ previous context
	\end{center}
	
	\vspace{0.3cm}
	
	\begin{alertblock}{Problem}
		The relevant information may appear far in the past.
	\end{alertblock}
	
\end{frame}

\begin{frame}
	\frametitle{\insertshortsubtitle: \insertsection}
	\framesubtitle{\insertsubsection: Some fun}
	
	\begin{center}
		\vgraphpage{humor/humor-words-order.jpg}
	\end{center}
	
\end{frame}

\subsection{Markov assumption}

\begin{frame}
	\frametitle{\insertshortsubtitle: \insertsection}
	\framesubtitle{\insertsubsection}
	
	\begin{block}{Sequence prediction}
		Given an input sequence:
		\[
		x = (x_1, x_2, \dots, x_n)
		\]
		predict:
		\[
		y = (y_1, y_2, \dots, y_n)
		\]
	\end{block}
	
	\vspace{0.3cm}
	
	Traditional approaches approximate:
	
	\[
	p(y|x)
	=
	\prod_i
	p(y_i \mid y_{i-l}, \dots, y_{i-1}, x_i)
	\]
	
	\vspace{0.3cm}
	
	\begin{alertblock}{Markov assumption}
		The current prediction depends only on a limited history window.
	\end{alertblock}
	
\end{frame}

\begin{frame}
	\frametitle{\insertshortsubtitle: \insertsection}
	\framesubtitle{\insertsubsection: Fixed-context window}
	
	\begin{center}
		
		\[
		\dots \quad y_{t-5} \quad y_{t-4} \quad y_{t-3}
		\quad
		\fbox{$y_{t-2}$}
		\quad
		\fbox{$y_{t-1}$}
		\quad
		y_t
		\]
		
	\end{center}
	
	\begin{exampleblock}{Example of a fixed-context window (size 3)}
		\textbf{Task}: Part-of-speech tagging
		
		\textbf{Input}: The animal did not cross the street because \textbf{it} was too tired.
		
		\textbf{Output}: DET NOUN VERB PREP VERB 
		\fbox{DET} \fbox{NOUN} \fbox{PREP} \textbf{??}
	\end{exampleblock}
	
	Only a limited portion of the past is accessible.
	
	\begin{block}{Consequence}
		Dependencies outside the context window cannot be modeled directly.
	\end{block}
	
\end{frame}

\begin{frame}
	\frametitle{\insertshortsubtitle: \insertsection}
	\framesubtitle{\insertsubsection: Some fun}
	
	\begin{center}
		\vgraphpage{humor/humor-markov.jpg}
	\end{center}
	
\end{frame}

\subsection{Limitations}

\begin{frame}
	\frametitle{\insertshortsubtitle: \insertsection}
	\framesubtitle{\insertsubsection: Loss of long-term context}
	
	\begin{exampleblock}{Long dependency example}
		``I grew up in \textbf{Algeria} \dots I know what \textbf{Sahhit} means.''
	\end{exampleblock}
	
%	\vspace{0.3cm}
	
	The relation between: Algeria $\rightarrow$ Sahhit may span many intermediate words.
	
%	\vspace{0.3cm}
	
	\begin{alertblock}{Problem}
		A fixed-size context window may miss relevant information.
	\end{alertblock}
	
\end{frame}

\begin{frame}
	\frametitle{\insertshortsubtitle: \insertsection}
	\framesubtitle{\insertsubsection: Need for memory}
	
	\begin{block}{Desired properties}
		A sequence model should:
		\begin{itemize}
			\item Process variable-length sequences
			\item Preserve past information
			\item Learn contextual dependencies
			\item Share parameters across time
		\end{itemize}
	\end{block}

	\begin{block}{Idea}
		\centering
		Introduce a hidden state acting as a dynamic memory.
		
		$\Longrightarrow$ 
		
		Recurrent Neural Networks
	\end{block}
	
\end{frame}



\section{Recurrent neural networks}

\begin{frame}
	\frametitle{\insertshortsubtitle}
	\framesubtitle{\insertsection}
	
	\begin{block}{From feedforward to sequential modeling}
		Standard neural networks assume independent inputs:
		\[
		y = f(x)
		\]
		
		They cannot naturally model temporal structure.
	\end{block}
	
	\begin{block}{Key idea}
		Introduce a state that evolves over time:
		\[
		h_t = \text{memory of past inputs}
		\]
		
		\[
		\Rightarrow \text{Recurrent Neural Networks}
		\]
	\end{block}
	
\end{frame}

\begin{frame}
	\frametitle{\insertshortsubtitle}
	\framesubtitle{\insertsection}
	
	\begin{center}
		\vgraphpage{humor/humor-rnn-memory.jpg}
	\end{center}
	
\end{frame}



\subsection{Architecture}

\begin{frame}
	\frametitle{\insertshortsubtitle: \insertsection}
	\framesubtitle{\insertsubsection}
	
	\begin{minipage}{0.49\textwidth} 
		\begin{itemize}
			\item \optword{Elman network}
			\begin{align*}
				h_t = f(x_t W_x + \textcolor{red}{h_{t-1}} W_h + b_h) \\
				\hat{y}_t = g(h_t W_y + b_y)
			\end{align*}
			
			\item \optword{Jordan network}
			\begin{align*}
				h_t = f(x_t W_x + \textcolor{red}{\hat{y}_{t-1}} W_h + b_h) \\
				\hat{y}_t = g(h_t W_y + b_y)
			\end{align*}
			
			\item The same parameters are reused at all time steps.
		\end{itemize}
	\end{minipage}
	%
	\begin{minipage}{0.5\textwidth}
		\hgraphpage[\textwidth]{RNN.pdf}
	\end{minipage}
	
\end{frame}

\begin{frame}
	\frametitle{\insertshortsubtitle: \insertsection}
	\framesubtitle{\insertsubsection: Implementation}
	
%	\begin{center}
%		\hgraphpage[0.65\textwidth]{RNN_impl.pdf}
%	\end{center}
	
	\begin{minipage}{0.65\textwidth}
		\hgraphpage[\textwidth]{RNN_impl.pdf}
	\end{minipage}
	\begin{minipage}{0.3\textwidth}
%		\scriptsize	
		\tiny
		$X \in \mathbb{R}^{m \times t \times 2} $
		
		\vskip5pt
		$ Y \in \mathbb{R}^{m \times t \times 4}$
		
		\vskip5pt
		$ H \in \mathbb{R}^{m \times t \times 3}$
		
		\vskip5pt
		$W_x \in \mathbb{R}^{2 \times 3} = 
		\begin{bmatrix}
			0.5 & -0.5 & -2. \\
			1. & -2.5 & 0.\\
		\end{bmatrix}$
		
		\vskip5pt
		$W_h \in \mathbb{R}^{3 \times 3}  = 
		\begin{bmatrix}
			-0.5 & 0. & -0.5\\
			0. & 0.5 & 0.5 \\
			1. & 0. & -0.5 \\
		\end{bmatrix}$
		
		\vskip5pt
		$b_h \in \mathbb{R}^{3} = [0., 2., 2.5]$
		
		\vskip5pt
		$W_y \in \mathbb{R}^{3 \times 4} = $
		
		$\begin{bmatrix}
			-1. & -1.5 & -1.5 & 1. \\
			2.  & -1.5 & -2.5 & 0. \\
			0.  & -1.5 & 1.5  & -1.5\\ 
		\end{bmatrix}$
		
		\vskip5pt
		$b_y \in \mathbb{R}^{4} = [-3., 4., -2.5, -4.]$
	\end{minipage}
	
\end{frame}

\subsection{Forward propagation and applications}

\begin{frame}
	\frametitle{\insertshortsubtitle: \insertsection}
	\framesubtitle{\insertsubsection: Forward propagation}
	
	Unfolding through time:
	\begin{itemize}
		\item The recurrent connection creates temporal dependencies.
		\item The same cell is reused at every time step.
		\item Unfolding converts the recurrent graph into a deep feed-forward graph.
	\end{itemize}
	
	\hgraphpage[\textwidth]{RNN-exp.pdf}
	
\end{frame}


\begin{frame}
	\frametitle{\insertshortsubtitle: \insertsection}
	\framesubtitle{\insertsubsection: Applications}
	
	\begin{tabular}{p{.32\textwidth}p{.15\textwidth}p{.4\textwidth}}
		%					\hline\hline
		\textbf{Type} & \textbf{Illustration} & \textbf{Example} \\
		%					\hline
		Many to many & 
		\vgraphpage[1.4cm, valign=c]{RNN_m2m1.pdf} & 
		Named entity detection \\
		
		%					\hline
		Many to many (Seq2seq) & 
		\vgraphpage[1.4cm, valign=c]{RNN_m2m2.pdf} & 
		Machine translation \\
		
		%					\hline
		Many to one & 
		\vgraphpage[1.4cm, valign=c]{RNN_m2o.pdf} & 
		Sentiment classification \\
		
		%					\hline
		One to many & 
		\vgraphpage[1.4cm, valign=c]{RNN_o2m.pdf} & 
		Image caption generation \\
		
		%					\hline\hline
		
	\end{tabular}
	
\end{frame}

\begin{frame}
	\frametitle{\insertshortsubtitle: \insertsection}
	\framesubtitle{\insertsubsection: Applications (Synchronous many-to-many)}
	
	\begin{minipage}{0.53\textwidth}
		\begin{itemize}
			\item Bias terms are omitted for simplicity.
			\item $h_0$ is a predefined victor (zeroes mostly).
		\end{itemize}
		
		\[ h_t = f(x_t, h_{t-1}, W_x, W_h)\]
		
		\[ \hat{y}_t = g(h_t, W_y)\]
		
		\[ J(\hat{y}, y) = \frac{1}{T} \sum_{t=1}^{T} j(\hat{y}_t, y_t)\]
	\end{minipage}
	\begin{minipage}{0.45\textwidth}
		\hgraphpage[\textwidth]{RNN_m2m1_exp.pdf}
	\end{minipage}
	
\end{frame}

\begin{frame}
	\frametitle{\insertshortsubtitle: \insertsection}
	\framesubtitle{\insertsubsection: Applications (Seq2seq)}
	
	\begin{minipage}{0.45\textwidth}
		\centering\vskip8pt
		\textbf{Encoder}
		
		$h_0$ is a predefined victor
		
		\[ h^e_t = f^e(x_t, h^e_{t-1}, W^e_x, W^e_h) \]
		
		\[ J(\hat{y}, y) = \frac{1}{T'} \sum_{t=1}^{T'} j(\hat{y}_{t}, y_t)\]
		
	\end{minipage}
	\begin{minipage}{0.45\textwidth}
		\centering
		\textbf{Decoder}
		
		$h^d_{0} = h^e_{T}$, $y_0$ is a predefined victor
		
		\[h^d_{t} = f^d(\hat{y}_{t-1}, h^d_{t-1}, W^d_x, W^d_h)\]
		
		\[ \hat{y}_{t} = g^d(h^d_{t-1}, W^d_y)\]
	\end{minipage}
	
	\vspace{-12pt}
	\begin{center}
		\hgraphpage[0.8\textwidth]{RNN_m2m2_exp.pdf}
	\end{center}
	
\end{frame}

\begin{frame}
	\frametitle{\insertshortsubtitle: \insertsection}
	\framesubtitle{\insertsubsection: Applications (Many-to-one)}
	
	\begin{minipage}{0.4\textwidth}
		\begin{center}
			$h_0$ is a predefined victor
		\end{center}
		
		\[ h_t = f(x_t, h_{t-1}, W_x, W_h)\]
		
		\[ \hat{y} = g(h_T, W_y)\]
		
		\[ J(\hat{y}, y) = j(\hat{y}, y)\]	
	\end{minipage}
	\begin{minipage}{0.58\textwidth}
		\hgraphpage{RNN_m2o_exp.pdf}
	\end{minipage}
	
\end{frame}

\begin{frame}
	\frametitle{\insertshortsubtitle: \insertsection}
	\framesubtitle{\insertsubsection: Applications (One-to-many)}
	
	\begin{minipage}{0.4\textwidth}
		
		\[h_1  = f(y_{0}, x, W_x, W_h)\]
		
		\[h_{t} = f(\hat{y}_{t-1}, h_{t-1}, W_y, W_h)\]
		
		\[\hat{y}_{t} = g(h_{t}, W_y)\]
		
		\[ J(\hat{y}, y) = \frac{1}{T} \sum_{t=1}^{T} j(\hat{y}_{t}, y_t)\]
	\end{minipage}
	\begin{minipage}{0.58\textwidth}
		\hgraphpage{RNN_o2m_exp.pdf}
	\end{minipage}
	
\end{frame}

\subsection{Backpropagation through time}

\begin{frame}
	\frametitle{\insertshortsubtitle: \insertsection}
	\framesubtitle{\insertsubsection}
	
	\begin{itemize}
		
		\item Define a cost function for each output.
		
		\item Use backpropagation through time (BPTT) \cite{1990-werbos}
		\begin{itemize}
			
			\item Unroll the recurrent network over time.
			
			\item The resulting architecture becomes similar to a feed-forward network.
			
			\item Compute gradients backward starting from the last time step.
			
			\item Accumulate gradients across all time steps.
			
			\item Update the parameters after processing the whole sequence.
			
		\end{itemize}
		
	\end{itemize}
	
\end{frame}

\begin{frame}
	\frametitle{\insertshortsubtitle: \insertsection}
	\framesubtitle{\insertsubsection: Synchronous many-to-many}
	
	\begin{align*}
		\frac{\partial J}{\partial W_y} & =
		\frac{1}{T} \sum_{t=1}^{T} \frac{\partial j(\hat{y}_t, y_t)}{\partial W_y}
		= \frac{1}{T} \sum_{t=1}^{T} \frac{\partial j(\hat{y}_t, y_t)}{\partial \hat{y}_t} 
		\frac{\partial g(h_t, W_y)}{\partial W_y} \\
		%		
		\frac{\partial J}{\partial W_h} & = \frac{1}{T} \sum_{t=1}^{T} \frac{\partial j(\hat{y}_t, y_t)}{\partial W_h}
		= \frac{1}{T} \sum_{t=1}^{T} \frac{\partial j(\hat{y}_t, y_t)}{\partial \hat{y}_t} 
		\frac{\partial g(h_t, W_y)}{\partial h_t} 
		\frac{\partial h_t}{\partial W_h} \\
		\frac{\partial h_t}{\partial W_h} & = 
		\frac{\partial f(x_t, h_{t-1}, W_x, W_h)}{\partial W_h} + 
		\frac{\partial f(x_t, h_{t-1}, W_x, W_h)}{\partial h_{t-1}} \frac{\partial h_{t-1}}{\partial W_h} \\
		%		
		\frac{\partial J}{\partial W_x} & = 
		\frac{1}{T} \sum_{t=1}^{T} \frac{\partial j(\hat{y}_t, y_t)}{\partial W_x}
		= \frac{1}{T} \sum_{t=1}^{T} \frac{\partial j(\hat{y}_t, y_t)}{\partial \hat{y}_t} 
		\frac{\partial g(h_t, W_y)}{\partial h_t} 
		\frac{\partial h_t}{\partial W_x} \\
		\frac{\partial h_t}{\partial W_x} & = 
		\frac{\partial f(x_t, h_{t-1}, W_x, W_h)}{\partial W_x} + 
		\frac{\partial f(x_t, h_{t-1}, W_x, W_h)}{\partial h_{t-1}} \frac{\partial h_{t-1}}{\partial W_x} \\
	\end{align*}
	
\end{frame}

\begin{frame}
	\frametitle{\insertshortsubtitle: \insertsection}
	\framesubtitle{\insertsubsection: Seq2seq}
	
	\begin{itemize}
		\item Decoder cell is updated as synchronous many-to-many, replacing $x_t$ by $\hat{y}_{t-1}$
		\item Encoder cell is updated as synchronous many-to-many, but without direct output
	\end{itemize}
	
	\[\frac{\partial J}{\partial W^e_h} = \frac{\partial J}{\partial h^e_T} \frac{\partial h^e_T}{\partial W^e_h} 
	\ \ \ \
	\frac{\partial J}{\partial W^e_x}  = \frac{\partial J}{\partial h^e_T} \frac{\partial h^e_T}{\partial W^e_x} 
	\]
	\begin{align*}
		%		\frac{\partial J}{\partial W^e_h} & = \frac{\partial J}{\partial h^e_t} \frac{\partial h^e_t}{\partial W^e_h} 
		%		= \frac{\partial J}{\partial h^e_T} \prod_{t=1}^{T} \frac{\partial h^e_t}{\partial W^e_h}  \\
		%%
		%		\frac{\partial J}{\partial W^e_x} & = \frac{\partial J}{\partial h^e_t} \frac{\partial h^e_t}{\partial W^e_x} 
		%		= \frac{\partial J}{\partial h^e_T} \prod_{t=1}^{T} \frac{\partial h^e_t}{\partial W^e_x}  \\
		%		
		\frac{\partial J}{\partial h^e_T} & = \frac{\partial J}{\partial h^d_0}
		= \frac{1}{T'} \sum_{t=1}^{T'} \frac{\partial j(\hat{y}_t, y_t)}{\partial h^d_0} 
		= \frac{1}{T'} \sum_{t=1}^{T'} 
		\frac{\partial j(\hat{y}_t, y_t)}{\partial \hat{y}_t} 
		\frac{\partial g(h^d_t, W^d_y)}{\partial h^d_t} 
		\frac{\partial h^d_t}{\partial h^d_0}\\
		%		 
		\frac{\partial h^d_t}{\partial h^d_0} & = 
		\frac{\partial f^d(\hat{y}_{t-1}, h^d_{t-1}, W^d_x, W^d_h)}{\partial h^d_{t-1}} \frac{\partial h^d_{t-1}}{\partial h^d_{0}} 
		= \prod_{i=1}^{t} \frac{\partial f^d(\hat{y}_{i-1}, h^d_{i-1}, W^d_x, W^d_h)}{\partial h^d_{i-1}}\\
	\end{align*}
	
\end{frame}

\begin{frame}
	\frametitle{\insertshortsubtitle: \insertsection}
	\framesubtitle{\insertsubsection: Many-to-one}
	
	\begin{align*}
		\frac{\partial J}{\partial W_y} & = \frac{\partial j(\hat{y}, y)}{\partial W_y}
		= \frac{\partial j(\hat{y}, y)}{\partial \hat{y}} \frac{\partial g(h_T, W_y)}{\partial W_y} \\
		%		
		\frac{\partial J}{\partial W_h} & = \frac{\partial j(\hat{y}, y)}{\partial W_h}
		= \frac{\partial j(\hat{y}, y)}{\partial \hat{y}} 
		\frac{\partial g(h_T, W_y)}{\partial h_T} 
		\frac{\partial h_T}{\partial W_h} \\
		\frac{\partial h_t}{\partial W_h} & = 
		\frac{\partial f(x_t, h_{t-1}, W_x, W_h)}{\partial W_h} + 
		\frac{\partial f(x_t, h_{t-1}, W_x, W_h)}{\partial h_{t-1}} \frac{\partial h_{t-1}}{\partial W_h} \\
		%		
		\frac{\partial J}{\partial W_x} & = \frac{\partial j(\hat{y}, y)}{\partial W_x}
		= \frac{\partial j(\hat{y}, y)}{\partial \hat{y}} 
		\frac{\partial g(h_T, W_y)}{\partial h_T} 
		\frac{\partial h_T}{\partial W_x} \\
		\frac{\partial h_t}{\partial W_x} & = 
		\frac{\partial f(x_t, h_{t-1}, W_x, W_h)}{\partial W_x} + 
		\frac{\partial f(x_t, h_{t-1}, W_x, W_h)}{\partial h_{t-1}} \frac{\partial h_{t-1}}{\partial W_x} \\
	\end{align*}
	
\end{frame}

\begin{frame}
	\frametitle{\insertshortsubtitle: \insertsection}
	\framesubtitle{\insertsubsection: One-to-many}
	
	\begin{itemize}
		
		\item The model generates a sequence from a single input.
		
		\item The initial hidden state is obtained from the input representation $h_0 = \phi(x)$;
		where $\phi$ is an encoder (e.g., CNN for images).
		
		\item The decoder receives sequential inputs:
		\begin{itemize}
			\item a special start token at $t=1$,
			\item previous target tokens during training,
			\item previous predictions during inference.
		\end{itemize}
		
		\item The recurrent decoder is trained similarly to synchronous many-to-many.
		
		\item Gradients are accumulated through time using BPTT.
		
	\end{itemize}
	
\end{frame}

\section{Limitations}

\begin{frame}
	\frametitle{\insertshortsubtitle}
	\framesubtitle{\insertsection}
	
	\begin{alertblock}{Main limitations of standard RNNs}
		\begin{itemize}
			\item Difficulty learning long-term dependencies
			\item Vanishing gradients
			\item Exploding gradients
			\item Sequential computation
		\end{itemize}
	\end{alertblock}
	
	\vspace{0.4cm}
	
	\begin{center}
		These limitations motivated gated architectures such as LSTM and GRU.
	\end{center}
	
\end{frame}

\subsection{Vanishing and exploding gradients}

\begin{frame}
	\frametitle{\insertshortsubtitle: \insertsection}
	\framesubtitle{\insertsubsection: Gradient flow}
	
	\begin{itemize}
		\item During BPTT, gradients propagate through all previous states.
		\item This introduces repeated multiplications:
	\end{itemize}
	
	\vspace{0.2cm}
	
	\[
	\frac{\partial h_t}{\partial h_{t-1}}
	\frac{\partial h_{t-1}}{\partial h_{t-2}}
	\cdots
	\frac{\partial h_1}{\partial h_0}
	\]
	
	\vspace{0.3cm}
	
	\begin{block}{Problem}
		Repeated multiplication may:
		\begin{itemize}
			\item shrink gradients toward zero,
			\item or make them grow uncontrollably.
		\end{itemize}
	\end{block}
	
\end{frame}

\begin{frame}
	\frametitle{\insertshortsubtitle: \insertsection}
	\framesubtitle{\insertsubsection: Formulas}
	
	$\frac{\partial h_t}{\partial W_h}$ can be simplified:
	\begin{align*}
		\frac{\partial h_t}{\partial W_h}  = &
		\frac{\partial f(x_t, h_{t-1}, W_x, W_h)}{\partial W_h} \\
		& + 
		\sum_{i=1}^{t-1} \left( {\color{red} \prod_{j=i+1}^{t} \frac{\partial f(x_j, h_{j-1}, W_x, W_h)}{\partial h_{j-1}}} \right)
		\frac{\partial f(x_i, h_{i-1}, W_x, W_h)}{\partial W_h}\\
	\end{align*}
	
	\begin{itemize}
		\item Can result in \keyword{vanishing gradients} (e.g. \expword{$f=\sigma(\sum)$})
		\item Also, \keyword{exploding gradients} (e.g. \expword{$f=ReLU(\sum)$})
		\item Possible solution: Truncating time steps \cite{2002-jaeger}
	\end{itemize}
	
\end{frame}

\subsection{Long-term dependency problem}

\begin{frame}
	\frametitle{\insertshortsubtitle: \insertsection}
	\framesubtitle{\insertsubsection}
	
	\begin{exampleblock}{Example}
		``I grew up in \textbf{Algeria} \dots I know what \textbf{Sahhit} means.''
	\end{exampleblock}
	
	The model must retain relevant information across many time steps.
	
	\begin{alertblock}{Core limitation}
		The hidden state $h_t \in \mathbb{R}^H$ is a fixed-size summary of all past inputs.
	\end{alertblock}
	
	\begin{itemize}
		\item Information must be continuously compressed at every step
		\item Old information can be overwritten by new inputs
		\item Long-range dependencies become harder to retrieve
	\end{itemize}
	
	\begin{center}
		This is a representational bottleneck, independent of gradient issues.
	\end{center}
	
\end{frame}

\subsection{Computational challenges}

\begin{frame}
	\frametitle{\insertshortsubtitle: \insertsection}
	\framesubtitle{\insertsubsection}
	
	\begin{block}{Sequential dependency}
		Each hidden state depends on the previous one:
		\[
		h_t = f(x_t, h_{t-1})
		\]
	\end{block}
	
	\begin{itemize}
		\item Computation across time steps is inherently sequential.
		\item Parallelism is possible across the batch dimension.
		\item Long sequences reduce hardware utilization efficiency.
		\item Memory usage grows with sequence length due to stored activations.
	\end{itemize}
	
	\begin{alertblock}{Consequence}
		RNNs are difficult to scale efficiently to very long sequences.
	\end{alertblock}
	
\end{frame}


\section{Advanced RNNs}

\begin{frame}
	\frametitle{\insertshortsubtitle}
	\framesubtitle{\insertsection}
	
	\begin{alertblock}{Why advanced recurrent architectures?}
		Standard RNNs suffer from:
		\begin{itemize}
			\item Vanishing gradients
			\item Difficulty modeling long-term dependencies
			\item Unstable training
		\end{itemize}
	\end{alertblock}
	
	\vspace{0.4cm}
	
	\begin{block}{Main idea}
		Introduce gating mechanisms to control information flow and memory updates.
	\end{block}
	
\end{frame}

\subsection{Long Short-Term Memory (LSTM)}

\begin{frame}
	\frametitle{\insertshortsubtitle: \insertsection}
	\framesubtitle{\insertsubsection: Architecture}
	
	\begin{minipage}{0.50\textwidth} 
		\begin{itemize}
			
			\item $i$: \optword{Input gate}: 
			Controls how much new information is added to the cell state.
			
			\item $f$: \optword{Forget gate}: 
			Controls how much information from the previous cell state is retained.
			
		\end{itemize}
	\end{minipage}
	%
	\begin{minipage}{0.49\textwidth}
		\hgraphpage[\textwidth]{LSTM.pdf}
	\end{minipage}
	
%	\vspace{0.2cm}
	
	\begin{itemize}
		\item $o$: \optword{Output gate}: 
		Controls how much information from the cell state is exposed to the hidden state.
		\item $g$: \optword{Candidate state}: 
		Represents new candidate information that may be added to the cell state.
	\end{itemize}
	
\end{frame}

\begin{frame}
	\frametitle{\insertshortsubtitle: \insertsection}
	\framesubtitle{\insertsubsection: Equations}
	\huge\vskip-24pt
	\begin{align*}
		f_t &= \sigma(W_f X_t + U_f h_{t-1} + b_f) \\
		i_t &= \sigma(W_i X_t + U_i h_{t-1} + b_i) \\
		g_t &= \tanh(W_g X_t + U_g h_{t-1} + b_g) \\
		o_t &= \sigma(W_o X_t + U_o h_{t-1} + b_o) \\
		c_t &= f_t \circ c_{t-1} \oplus i_t \circ g_t \\
		h_t &= o_t \circ \tanh(c_t)
	\end{align*}
	
\end{frame}

\begin{frame}
	\frametitle{\insertshortsubtitle: \insertsection}
	\framesubtitle{\insertsubsection: Gradient propagation}
	
	The cell state evolves as: $c_t = f_t \circ c_{t-1} + i_t \circ g_t$
	
	Therefore: \[\frac{\partial c_t}{\partial c_{t-1}} = f_t\]
	
	Over multiple time steps: \[\frac{\partial c_t}{\partial c_k} = \prod_{i=k+1}^{t} f_i\]
	
	\begin{itemize}
		\item Forget gates can preserve information: $f_t \approx 1$
		
		\item Gradients can flow through the cell state over long sequences.
		
		\item LSTMs mitigate the vanishing gradient problem.
	\end{itemize}
	
\end{frame}

\begin{frame}
	\frametitle{\insertshortsubtitle: \insertsection}
	\framesubtitle{\insertsubsection: Some humor}
	
	\begin{center}
		\hgraphpage[.5\textwidth]{humor/humor-lstm.jpg}
	\end{center}
	
\end{frame}

\subsection{Gated Recurrent Unit (GRU)}


\begin{frame}
	\frametitle{\insertshortsubtitle: \insertsection}
	\framesubtitle{\insertsubsection: Architecture}
	
	\begin{minipage}{0.50\textwidth} 
		\begin{itemize}
			\item $z$: \optword{Update gate};
			Merges LSTM's input and forget gates. 
			\item $r$: \optword{Reset gate}; 
			How much update from the past state?
			\item $h'$: \optword{Candidate hidden state};
			Current state.	
		\end{itemize}
	\end{minipage}
	%
	\begin{minipage}{0.49\textwidth}
		\hgraphpage[\textwidth]{GRU.pdf}
	\end{minipage}
	
\end{frame}

\begin{frame}
	\frametitle{\insertshortsubtitle: \insertsection}
	\framesubtitle{\insertsubsection: Equations}
	
	{\huge\vskip-24pt
		\begin{align*}
			z_t &= \sigma(W_z X_t + U_z h_{t-1} + b_z) \\
			r_t &= \sigma(W_r X_t + U_r h_{t-1} + b_r) \\
			h'_t &= \tanh(W X_t + U (h_{t-1} \circ r_t) + b) \\
			h_t &= z_t \circ h_{t-1} \oplus (1-z_t) \circ h'_t
		\end{align*}
	}
	
\end{frame}

\begin{frame}
	\frametitle{\insertshortsubtitle: \insertsection}
	\framesubtitle{\insertsubsection: Gradient propagation}
	
	The hidden state evolves as:
	
	\[
	h_t = z_t \circ h_{t-1} + (1-z_t) \circ h'_t
	\]
	
	\vspace{0.3cm}
	
	Therefore:
	
	\[
	\frac{\partial h_t}{\partial h_{t-1}}
	=
	z_t
	+
	(1-z_t)
	\frac{\partial h'_t}{\partial h_{t-1}}
	\]
	
	\vspace{0.3cm}
	
	\begin{itemize}
		\item The update gate can preserve past information:
		\[
		z_t \approx 1
		\]
		
		\item A direct path is maintained between successive hidden states.
		
		\item GRUs facilitate gradient propagation across long sequences.
	\end{itemize}
	
\end{frame}

\begin{frame}
	\frametitle{\insertshortsubtitle: \insertsection}
	\framesubtitle{\insertsubsection: Some humor}
	
	\begin{center}
		\vgraphpage{humor/humor-gru.png}
	\end{center}
	
\end{frame}

\subsection{Bidirectional and deep RNNs}


\begin{frame}
	\frametitle{\insertshortsubtitle: \insertsection}
	\framesubtitle{\insertsubsection: Bidirectional RNN}
	
	\begin{center}
		\hgraphpage[0.8\textwidth]{bi-rnn_exp.pdf}
	\end{center}
	
\end{frame}

\begin{frame}
	\frametitle{\insertshortsubtitle: \insertsection}
	\framesubtitle{\insertsubsection: Deep RNN}
	
	\begin{center}
		\hgraphpage[0.8\textwidth]{deep-rnn_exp.pdf}
	\end{center}
	
\end{frame}

%\subsection{Applications in NLP and time series}
%
%\begin{frame}
%	\frametitle{\insertshortsubtitle: \insertsection}
%	\framesubtitle{\insertsubsection}
%	
%\end{frame}

\insertbibliography{DL03L-RNN}{*}

\end{document}